
\section{D\'efinitions des notions fondamentales}\label{sec:definitions}

\subsection{La notion de comp\'etence}
% - Définitions classiques
% - Approches pédagogiques
% - Dimension évaluative

La notion de comp\'etence est l'une des plus discut\'ees en sciences de l'\'education. 
Elle n'appartient pas originellement au monde scolaire: elle est importée du monde du travail et de la formation professionnelle, o\`u elle d\'esigne la capacit\'e op\'erationnelle \`a accomplir des t\^aches dans des contextes d\'etermin\'es. 
Son introduction dans le cadre scolaire, amorc\'ee dans les ann\'ees 1990, a suscit\'e de vives controverses, certains y lisant une r\'eduction utilitariste de l'id\'eal \'educatif.

Une d\'efinition de r\'ef\'erence en p\'edagogie francophone est celle de~\cite{perrenoud2011construire}:

\begin{displayquote}[\cite{perrenoud2011construire}]
Une capacit\'e d'action efficace face \`a une famille de situations, qu'on arrive \`a ma\^\i triser parce qu'on dispose \`a la fois des connaissances n\'ecessaires et de la capacit\'e de les mobiliser \`a bon escient, en temps opportun, pour identifier et r\'esoudre de vrais probl\`emes.
\end{displayquote}


Deux \'el\'ements ressortent de cette d\'efinition. 
Premi\`erement, la comp\'etence n'est pas la simple possession de savoirs: elle requiert leur \emph{mobilisation} en situation.
Deuxi\`emement, elle s'exerce face \`a des \og{}familles de situations\fg{}, c'est-\`a-dire qu'elle suppose une capacit\'e de transfert, d'adaptation \`a des contextes in\'edits.

\cite{tardif2006evaluation}compl\`ete cette d\'efinition en insistant sur la mobilisation de \og{}ressources internes et externes\fg{}. % chktex 2
Cette pr\'ecision est très importante: elle ouvre la voie \`a la reconnaissance que les outils et les bases de donn\'ees font partie d'outils légitimes à disposition. 
C'est pr\'ecis\'ement cette ouverture qui rend la question des IAG si d\'elicate: si la comp\'etence inclut la mobilisation de ressources externes, alors l'IA g\'en\'erative peut-elle \^etre consid\'er\'ee comme une telle ressource? S'il existe un outil capable de répondre à une famille de situations, est-ce que la comp\'etence consiste à savoir utiliser cet outil, ou à faire sans lui?

Au plan institutionnel, la France a int\'egr\'e l'approche par comp\'etences dans ses curricula via le \emph{Socle commun de connaissances, de comp\'etences et de culture} (2006, r\'evis\'e en 2015), organisant les apprentissages autour de domaines transversaux. 
Cette architecture curriculaire vise \`a former des \'el\`eves capables de mobiliser leurs savoirs dans des situations de vie r\'eelle --- objectif qui se heurte de plein fouet aux possibilit\'es des IAG de simuler cette mobilisation~\citep{leroux2015evaluer}.

\subsection{La notion de r\'eussite scolaire}
% - Indicateurs traditionnels (notes, diplômes)
% - Vision institutionnelle
% - Limites conceptuelles

Dans son acception la plus courante, la r\'eussite scolaire renvoie \`a l'obtention de notes satisfaisantes, au passage dans la classe sup\'erieure, \`a la certification par un dipl\^ome.
C'est une r\'eussite au sens \emph{institutionnel}: est r\'eussissant l'\'el\`eve que l'institution reconna\^\i t comme tel, \`a travers ses instruments de mesure et de classement.
Cette d\'efinition, aussi fonctionnelle qu'elle soit, masque une tension fondamentale: ce qu'elle mesure rel\`eve-t-il de comp\'etences authentiquement acquises, ou d'une conformit\'e aux codes et attendus de l'institution~?

La sociologie critique a depuis longtemps soulign\'e le caract\`ere socialement construit de la r\'eussite scolaire. 
Pour Bourdieu et Passeron, l'\'ecole ne r\'ecompense pas le m\'erite individuel: elle valorise le \emph{capital culturel} que les familles transmettent in\'egalement selon leur position sociale~\citep{pluet1971pierre}. 
La r\'eussite scolaire appara\^\i t d\`es lors comme une performance sociale autant que cognitive: les \'el\`eves issus de milieux favoris\'es ma\^\i trisent plus rapidement les codes que l'institution r\'ecompense. 
L'\'ecole l\'egitime ainsi des in\'egalit\'es sociales de d\'epart en les faisant passer pour des diff\'erences de m\'erite ou d'intelligence.

Ce constat sociologique est crucial pour notre sujet: si la r\'eussite scolaire \'etait d\'ej\`a partiellement d\'econnect\'ee de la comp\'etence r\'eelle, l'arriv\'ee des IAG ne fait qu'aggraver cette d\'econnexion. 
De plus, l'article de~\cite{alexandre2019ia} dans \emph{Pouvoirs} pose la question radicale: dans une \og{}\'economie de la connaissance\fg{} o\`u l'intelligence est la cl\'e de tous les pouvoirs, l'\'ecole, \og{}en tant que technologie de transmission de l'intelligence, est d'ores et d\'ej\`a une technologie d\'epass\'ee\fg{}~\citep{alexandre2019ia}.
Si cette th\`ese est peut \^etre vue comme excessive dans sa radicalit\'e, elle pointe un fait r\'eel: les indicateurs traditionnels de r\'eussite (notes, dipl\^omes) r\'ev\`elent de plus en plus leurs \emph{limites conceptuelles} face \`a un monde o\`u les outils cognitifs externalis\'es transforment profond\'ement ce que signifie \og{}savoir\fg{}.

\subsection{Les intelligences artificielles g\'en\'eratives}
% - Définition technique rapide
% - Exemples d'usages en contexte scolaire

Les intelligences artificielles g\'en\'eratives (IAG) d\'esignent des syst\`emes capables de produire, \`a partir d'une instruction en langage naturel, des contenus vari\'es: textes, images, code, synth\`eses, argumentations. 
Techniquement, elles reposent sur de grands mod\`eles de langage (LLM, \emph{Large Language Models}) entra\^\i n\'es sur des corpus massifs, qui g\'en\`erent des r\'eponses par pr\'ediction statistique de la suite la plus probable d'une s\'equence de tokens (ceci est une simplification mais elle suffit dans ce contexte). 
ChatGPT (OpenAI), Gemini (Google), Claude (Anthropic) ou Mistral illustrent ce paradigme.

Comme le retracent Jabraoui et Vandapuye (2024), l'usage de l'IA en \'education a une histoire de plus de soixante ans, des premiers syst\`emes experts des ann\'ees 1980 aux syst\`emes adaptatifs des ann\'ees 2000, avant l'arriv\'ee des IA g\'en\'eratives actuelles~\citep{siham2024intelligence}.
Ce qui diff\'erencie fondamentalement ces derni\`eres de leurs pr\'ed\'ecesseurs est leur \emph{g\'en\'eralit\'e}: elles ne sont pas sp\'ecialis\'ees dans un domaine ou une t\^ache, mais peuvent produire tout type de contenu textuel, y compris des types de productions scolaires qui avaient jusqu'alors servi d'\'etalon de la comp\'etence humaine.
Elles ne sont plus spécialistes d'une seule discipline (math\'ematiques, langues, etc.) ou d'une seule fonction (tuteur, correcteur, etc.), mais peuvent intervenir dans de multiples domaines et ainsi devenir des outils polyvalents répondant à une grande variété de besoins.

Dans le contexte scolaire, leurs usages sont multiples et difficiles \`a limiter: r\'edaction de dissertations et de rapports, g\'en\'eration de plans, r\'esolution d'exercices, aide \`a la compr\'ehension de textes complexes, simulation de dialogues, pr\'eparation aux examens. 
Contrairement \`a la calculatrice ou au correcteur orthographique --- outils \`a p\'erim\`etre d'action d\'elimit\'e qui avaient \'et\'e progressivement int\'egr\'es dans les pratiques scolaires --- les IAG op\`erent dans le domaine m\^eme que l'\'evaluation scolaire cherche \`a mesurer. 
Alors qu'une calculatrice peut aider \`a faire un calcul, on mesure chez l'\'etudiant sa capacit\'e \`a faire le calcul sans calculatrice mais aussi d'autres compétences math\'ematiques plus \'elabor\'ees, comme la mod\'elisation ou l'interpr\'etation de r\'esultats, chose qu'une calculatrice ne peut faire de par son p\'erim\`etre d'action limit\'e.
Cette superposition entre le domaine d'opération des IAG et le domaine de compétence évaluées dans le système scolaire est au c\oe{}ur de la probl\'ematique qui nous occupe.
